\documentclass{article}

%%% Defines %%%%%%%%%%%%%%%%%%%%%%%%%%%%%%%%%%%%%%%%%%%%%%%%%%%%%%%%%%%%%%%%%%%%

\usepackage{graphicx}
\usepackage{tikz}
\usepackage{pgfplots}
\pgfplotsset{compat=1.18}
\usetikzlibrary{shapes.misc}
\usepackage[european,straightvoltages]{circuitikz}
\usepackage[ngerman]{babel} 
\usepackage[colorlinks]{hyperref}
\usepackage{caption, float, subcaption}
\usepackage{xcolor}
\usepackage{setspace}
\usepackage{mathtools, amssymb, ntheorem, amsmath, siunitx}
\sisetup{per-mode=fraction, separate-uncertainty=true,exponent-base=10,output-decimal-marker={,}}
\DeclareSIUnit\liter{l}
\usepackage{enumitem}
\usepackage{minted}

\usepackage{tikz-uml}

% maths
\newcommand{\abs}[1]{\left| #1 \right|}
\newcommand{\br}[1]{\left( #1 \right)}
\newcommand{\ubar}[1]{\mkern 1.5mu\underline{\mkern-1.5mu#1\mkern-1.5mu}\mkern 1.5mu}

% image
\newcommand{\img}[5]{
    \begin{figure} [#5]
    \centering
    \includegraphics[width=#2\linewidth]{#1}
    \caption{#3}
    \label{pic:#4}
    \end{figure}
}

% tikz
\newenvironment{gfx}[3]
{
    \newcommand{\gfxname}{#2}
    \newcommand{\gfxcaption}{#1}
    \begin{figure} [#3]
    \centering
    \begin{tikzpicture}
}
{
    \end{tikzpicture}
    \caption{\gfxcaption}
    \label{gfx:\gfxname}
    \end{figure}
}


% cicuitikz
\newenvironment{ckt}[3]
{
    \newcommand{\cktname}{#2}
    \newcommand{\cktcaption}{#1}
    \begin{figure} [#3]
    \centering
    \begin{circuitikz}
	\draw
}
{
	;
    \end{circuitikz}
    \caption{\cktcaption}
    \label{ckt:\cktname}
    \end{figure}
}

% Code
\newfloat{Code}{htbp}{loc}
\floatname{Code}{Quelltext}
\definecolor{LightGray}{gray}{0.9}
\newenvironment{code}[4]
{%
  \VerbatimEnvironment
  \begin{Code} [#4]
  \caption{#2}%
  \label{cod:#3}%
  \begin{minted}[frame=lines,framesep=2mm,baselinestretch=1.2,bgcolor=LightGray,fontsize=\footnotesize,style=emacs]{#1}%
}
{%
  \end{minted}%
  \vspace{-20pt}%
  \end{Code}%
}
\makeatletter
\AtBeginEnvironment{minted}{\dontdofcolorbox}
\def\dontdofcolorbox{\renewcommand\fcolorbox[4][]{##4}}
\makeatother
\newcommand{\listofcode}{
  \doublespacing
  \listof{Code}{Quelltextverzeichnis}
}

%%% Head %%%%%%%%%%%%%%%%%%%%%%%%%%%%%%%%%%%%%%%%%%%%%%%%%%%%%%%%%%%%%%%%%%%%%%%

\title{Title}
\author{Benjamin Brohs, Kevin Keheck, Justin Meng}
\date{Date}

\setlength{\parindent}{0pt}
\setlength{\parskip}{1em}

\begin{document}

%%% Title %%%%%%%%%%%%%%%%%%%%%%%%%%%%%%%%%%%%%%%%%%%%%%%%%%%%%%%%%%%%%%%%%%%%%%

\begin{titlepage}
  \centering
	\begin{tabular}{lcr}
		\includegraphics[width=0.35\textwidth]{fachbereich.png} & \hspace{0.195\textwidth} & \includegraphics[width=0.35\textwidth]{Q04_HTW_Berlin_Logo_quer_pos_FARBIG_RGB.jpg}\\
	\end{tabular}	
	\\[3cm]
	\Large
	Belegarbeit\\
	\vspace{2cm}
	\textbf{Minecraft-One-Week-Challenge Reloaded}\\
	\vspace{2cm}
	\begin{tabular}{ll} 
		Im Studiengang: & Computer Engineering \\		
	\end{tabular}	
	\\[3cm]
	\normalsize
	\begin{tabular}{ll}
	      \textbf{Erstellt von:} & Benjamin Brohs, Kevin Keheck, Justin Meng \\
        \textbf{Modul:} & Softwaretechnik \\
        \textbf{Semester:} & Sommersemester 2025 \\
	\textbf{Dozent:} & Thomas Baar
	\end{tabular}	
\end{titlepage}

\tableofcontents

\newpage

%%% Inhalt %%%%%%%%%%%%%%%%%%%%%%%%%%%%%%%%%%%%%%%%%%%%%%%%%%%%%%%%%%%%%%%%%%%

\section{Meilenstein 1} \label{sec:ms1}

Meilenstein 1 befasst sich mit der Beschreibung der Software, wie sie ohne die Änderung aufgebaut ist. Sowohl textuell, als auch per Modell.

\subsection{Informelle Beschreibung} \label{subsec:inf}

Grundlegend basiert diese Softftware auf der Idee des Spiels \href{https://www.minecraft.net}{Minecraft}. Bei der vorgestellen Software jedoch, handelt es sich um eine wesentlich ausgedünnte Version dessen. Der Sinn soll sein, dass ein Spieler sich in einer Welt, generiert aus Blöcken verschiedener Texturen. Dem Spieler ist es möglich, diese Welt zu manipulieren. In erster Linie handelt es sich dabei um das Löschen von Blöcken. Die generierte Welt ist groß genug, sodass sich der Spieler in ihr frei bewegen kann. Der Spieler hat dabei die möglichkeit zu Springen, zu Sprinten und sich verlangsamt zu bewegen. Dabei stößt der Spieler auf verschiedenste generierte Strukturen. Das Original Open-Source-Projekt ist zu finden unter \url{https://github.com/Hopson97/MineCraft-One-Week-Challenge}.

\subsection{Use-Case-Modell} \label{subsec:usecase}

\begin{tikzpicture}
  % Systemgrenze
  \begin{umlsystem}[x=4, y=0]{Minecraft-Klon}
    \umlusecase[x=0, y=0]{Laufen}
    \umlusecase[x=0, y=-1.5]{Springen}
    \umlusecase[x=0, y=-3]{Schleichen}
    \umlusecase[x=0, y=-4.5]{Fliegen}
    \umlusecase[x=0, y=-6]{Block abbauen}
    \umlusecase[x=0, y=-7.5]{Block platzieren}
    \umlusecase[x=0, y=-9]{Maussteuerung umschalten}
  \end{umlsystem}

  % Akteur
  \umlactor[x=-2, y=-4]{Spieler}

  % Assoziationen
  \umlassoc{Spieler}{usecase-1}
  \umlassoc{Spieler}{usecase-2}
  \umlassoc{Spieler}{usecase-3}
  \umlassoc{Spieler}{usecase-4}
  \umlassoc{Spieler}{usecase-5}
  \umlassoc{Spieler}{usecase-6}
  \umlassoc{Spieler}{usecase-7}
\end{tikzpicture}

\subsection{Domänenmodell} \label{subsec:domain}

%%% Verzeichnisse %%%%%%%%%%%%%%%%%%%%%%%%%%%%%%%%%%%%%%%%%%%%%%%%%%%%%%%%%%%%%%
\section{Anhang}

In diesem Abschnitt finden sich verschiedene Verzeichnisse der genutzten 
Ressourcen, sowie zusätzliche Anhänge von Ressourcen, die nicht direkt im 
Hauptteil des Dokuments enthalten sind.

\listoffigures
\listoftables
\listofcode

\iffalse
\section*{Arbeitsverteilung}

\subsection*{Benjamin Brohs}
\begin{mylist}
  \mylistentry{sec:intro}
\end{mylist}
\fi

%%% Anhänge %%%%%%%%%%%%%%%%%%%%%%%%%%%%%%%%%%%%%%%%%%%%%%%%%%%%%%%%%%%%%%%%%%%%

\iffalse
\begin{table}[h]
	\centering
	\begin{tabular}{|c|c|}
	\hline
	$U_{AB}$ & I \\
	\hline
	0   & 0 \\
	0,1 & ~ \\
	\hline
	\end{tabular}
	\caption{Messwerte nach 3.3.1}
	\label{tab:rx}
\end{table}
\fi

\iffalse
\begin{gfx}{Bodediagramm Phasengang Bandpass}{bodeband2}{}
    \begin{axis}
        [
        xlabel={$f$ [Hz] },
        ylabel={Phasenunterschied $\Delta \varphi_{12}$ [°] },
        grid=both,
        grid style={line width=.1pt, draw=gray!10},
        major grid style={line width=.2pt,draw=gray!50},
        xmin=20, xmax=100000,
        ymin=-90, ymax=90,
        width=11.89cm, height=6cm,
        every axis plot/.append style={very thick, mark=*},
        legend style={at={(0.5,-0.1)}, anchor=north, legend columns=-1},
        xtick={10, 100, 1000, 10000, 100000},
        ytick={-90,-75,-60,-45,-30,-15,0,15,30,45,60,75,90},
        xmode=log
        ]

        \addplot[
        scatter,
        only marks,
        point meta=explicit symbolic
        ] table[meta=label] {
            x y label
            100 74 1
            200 59 2
            341.49 44 10
            500 32 3
            1000 14 4
            2000 -1 5
            5000 -20 6
            10000 -39 7
            12847.92 -46 11
            20000 -59 8
            50000 -76 9
        };

        \addplot[
        no marks,
        thin,
        ] coordinates{
            (20, 86.51900995)
            (23.79683391, 85.86013154)
            (28.3144652, 85.07750427)
            (33.68973128, 84.1485388)
            (40.08544699, 83.04696273)
            (47.6953362, 81.74251804)
            (56.74989969, 80.20084926)
            (67.52339687, 78.38375782)
            (80.34215301, 76.25009007)
            (95.59444355, 73.7576336)
            (113.7422548, 70.86648447)
            (135.3352773, 67.54432971)
            (161.0275558, 63.7738094)
            (191.5972999, 59.56138864)
            (227.9704562, 54.94591252)
            (271.2487541, 50.00361115)
            (322.7430775, 44.84581114)
            (384.0131705, 39.60725845)
            (456.9148818, 34.42684488)
            (543.6563777, 29.42650423)
            (646.8650261, 24.69507184)
            (769.6669794, 20.2811201)
            (915.7818637, 16.1944633)
            (1089.635445, 12.41311546)
            (1296.493686, 8.892029277)
            (1542.622245, 5.571078989)
            (1835.476267, 2.381199607)
            (2183.926194, -0.751349392)
            (2598.526445, -3.902071813)
            (3091.835111, -7.145930748)
            (3678.794331, -10.554851)
            (4377.182883, -14.19424224)
            (5208.154703, -18.11787564)
            (6196.879622, -22.3608024)
            (7373.305756, -26.93078726)
            (8773.066621, -31.80001815)
            (10438.56046, -36.90028499)
            (12420.23448, -42.12536058)
            (14778.11285, -47.34270134)
            (17583.61485, -52.41271888)
            (20921.7181, -57.2098854)
            (24893.53254, -61.63886918)
            (29619.36296, -65.64164712)
            (35242.35304, -69.19586452)
            (41932.82109, -72.30765117)
            (49893.41894, -75.00259162)
            (59365.27018, -77.317439)
            (70635.27373, -79.29373889)
            (84044.79385, -80.97351325)
            (100000, -82.3966604)
        };

        \addplot [densely dotted, thin, no marks] coordinates {(341.485847695166,-90) (341.485847695166,90)};
        \node [below left] at (axis cs:341.485847695166,43.14655286) {$f_h$};
        
        \addplot [densely dotted, thin, no marks] coordinates {(2094.60836705993,-90) (2094.60836705993,90)};
        \node [below left] at (axis cs:2094.60836705993,0) {$f_0$};
        
        \addplot [densely dotted, thin, no marks] coordinates {(12847.9239797778,-90) (12847.9239797778,90)};
        \node [above right] at (axis cs:12847.9239797778,-43.14655286) {$f_t$};
    \end{axis}
\end{gfx}
\fi

%%%%%%%%%%%%%%%%%%%%%%%%%%%%%%%%%%%%%%%%%%%%%%%%%%%%%%%%%%%%%%%%%%%%%%%%%%%%%%%%

\end{document}